% related-work.tex -- Enterprise Node Orchestrator Paper (RU-V5)

\section{Related Work}

\subsection{Production ZK Node Architectures}

The architecture of ZK L2 nodes has converged toward a multi-component
pattern where transaction sequencing, batch formation, proof generation,
and L1 submission are handled by separate processes.

Polygon zkEVM~\cite{polygon2024architecture} decomposes its node into a
Sequencer (transaction ordering), Aggregator (proof orchestration), zkProver
(recursive STARK-to-SNARK proving), and Synchronizer (L1 state recovery). The
Aggregator delegates proof tasks to the zkProver and manages retries on
failure. Hermez 2.0~\cite{hermez2023deepdive} introduced the Proof of
Efficiency consensus mechanism, establishing the batch lifecycle pattern
(sequence, aggregate, verify) that subsequent systems adopted.

zkSync Era~\cite{zksync2025lifecycle, zksync2025server} separates execution
from proving: the Sequencer processes transactions and produces soft
confirmations (sub-second), while the Batcher groups blocks into batches for
proving on a separate timeline. Recovery is idempotent---the Batcher resumes
from the first uncommitted batch. This architecture explicitly decouples
user-facing latency from proof generation latency.

Scroll~\cite{scroll2023architecture, scroll2025euclid} introduces a
two-stage proving pipeline: blocks are grouped into chunks, chunks are
proved individually, and chunk proofs are aggregated into batch proofs. A
Coordinator randomly assigns chunk proofs to a pool of provers, providing
horizontal scalability. Failed proof tasks are reassigned to other provers.

\subsection{Event-Driven ZK Orchestration}

The push0 framework~\cite{push0_2025orchestration}, deployed on Zircuit
since March 2025, represents the state of the art in ZK proof orchestration.
push0 uses stateless dispatchers and collectors coordinated via NATS
JetStream, achieving P50 orchestration latency of 5.2--5.3~ms and P99 of
8.0--8.5~ms on a 3-node Kubernetes cluster. Orchestration overhead is
0.05--0.1\% of proving time. Crash recovery achieves 0 task loss for
dispatchers and 0.5~s mean time to recovery (MTTR) for collectors. The
system scales to 32 concurrent dispatchers at 99\% efficiency.

Table~\ref{tab:architecture_comparison} compares the architectural patterns
across production systems and our proposed design.

\begin{table*}[t]
\centering
\caption{Production ZK Node Architecture Comparison}
\label{tab:architecture_comparison}
\resizebox{\textwidth}{!}{%
\begin{tabular}{@{}llllll@{}}
\toprule
\textbf{System} & \textbf{Architecture} & \textbf{Proving Model} & \textbf{Batch Pipeline} & \textbf{Crash Recovery} & \textbf{Scale Target} \\
\midrule
Polygon zkEVM & Seq + Agg + zkProver & Recursive STARK$\to$SNARK & Seq forms, Agg proves & Sync replays from L1 & Public L2 \\
zkSync Era    & Seq + Batcher + Prover & Execution $\neq$ proving & Blocks$\to$Batches$\to$L1 & Idempotent resume & Public L2 \\
Scroll        & Seq + Coord + Pool    & Coord assigns to pool & Blocks$\to$Chunks$\to$Batches & Reassign failed tasks & Public L2 \\
push0/Zircuit & Dispatchers + Bus     & Stateless dispatchers & Ranges$\to$Recursive$\to$Groth16 & At-least-once replay & Public L2 \\
\textbf{Ours} & \textbf{Single node, 3 loops} & \textbf{Child process} & \textbf{Txs$\to$Batch$\to$Prove$\to$Submit} & \textbf{WAL + checkpoint} & \textbf{Enterprise} \\
\bottomrule
\end{tabular}}
\end{table*}

\subsection{State Machine Design Patterns}

Petrasch~\cite{petrasch2018statemachine} describes the transformation of
state machines to event-driven architectures in enterprise systems,
establishing the pattern of typed state enumerations with explicit transition
guards. The Saga Orchestrator pattern~\cite{saga2024orchestrator} applies
state machines to long-running transactions with compensating actions,
directly analogous to the ZK proof pipeline where a failed proving step
must roll back batch state. Event sourcing with the Decider
pattern~\cite{eventsourcing2024typescript} provides a command-to-event
model for TypeScript implementations.

\subsection{Crash Recovery and Persistence}

Write-ahead logging (WAL) is the foundational crash recovery mechanism for
database systems. Mohan et al.~\cite{mohan1992aries} established the ARIES
protocol with checkpoint-based recovery and log sequence number (LSN)
tracking. Modern implementations for Node.js
applications~\cite{wal2024architecture} and graceful shutdown
patterns~\cite{nodejs2025shutdown} ensure that in-flight operations are
persisted before process termination. Our prior work (RU-V4) demonstrated
WAL-backed crash recovery for ZK batch aggregation with sub-150~$\mu$s
write latency and zero transaction loss~\cite{baylina2019smt}.

\subsection{ZK Prover Performance}

Groth16~\cite{groth2016size} remains the standard for succinct ZK proofs
with constant-size verification. Prover performance varies by implementation:
snarkjs (WASM)~\cite{snarkjs} achieves 12,757~ms for 274K constraints on
commodity hardware, while rapidsnark (C++ with native field arithmetic)
offers 4--10$\times$ speedup~\cite{mopro2025comparison}. ICICLE-Snark
(GPU-accelerated) achieves 63$\times$ MSM improvement~\cite{icicle2025snark},
reducing proving time to 1--3~s for circuits in the 1--3M constraint range.
The ZK proving infrastructure landscape~\cite{zkcloud2024landscape} shows a
trend toward prover decentralization and hardware acceleration.

\subsection{HTTP Framework Selection}

For the enterprise node API layer, framework selection directly impacts
throughput. Empirical comparisons~\cite{fastify2025comparison} demonstrate
that Fastify achieves 2--4$\times$ higher request throughput than Express.js
with native JSON schema validation and lower memory overhead. Fastify's
WebSocket integration~\cite{fastify2025websocket} enables real-time event
streaming for enterprise client applications.
